\documentclass{ximera}

\begin{document}
	\author{Stitz-Zeager}
	\xmtitle{Exercises for More Trigonometric Identities}{}

\mfpicnumber{1} \opengraphsfile{ExercisesforMoreTrigonometricIdentities} % mfpic settings added 



\begin{question}
    
In Exercises \ref{evenoddfirst} - \ref{evenoddlast}, use the Even / Odd Identities to verify the identity.  Assume all quantities are defined.

\begin{problem}\label{evenoddfirst}
    $\sin(3\pi - 2\theta) = -\sin(2\theta - 3\pi)$
\end{problem}

\begin{problem}
    $\cos \left( -\frac{\pi}{4} - 5t \right) = \cos \left( 5t + \frac{\pi}{4} \right)$
\end{problem}

\begin{problem}
     $\tan(-x^{2} + 1) = -\tan(x^{2} - 1)$
\end{problem}

\begin{problem}
    $\csc(-\theta - 5) = -\csc(\theta + 5)$
\end{problem}

\begin{problem}
    $\sec(-6x) = \sec(6x)$
\end{problem}

\begin{problem}\label{evenoddlast}
     $\cot(9 - 7\theta) = -\cot(7\theta - 9)$
\end{problem}

\end{question}


\begin{question}
    
In Exercises \ref{sumdifffirst} - \ref{sumdifflast}, use the Sum and Difference Identities to find the exact value.  You may have need of the Quotient, Reciprocal or Even / Odd Identities as well.

\begin{problem} \label{cos75} \label{sumdifffirst} 
    $\cos(75^{\circ})$

\begin{solution}
$\cos(75^{\circ}) = \frac{\sqrt{6} - \sqrt{2}}{4} $
\end{solution}
\end{problem}

\begin{problem}
$\sec(165^{\circ})$

\begin{solution}
$\sec(165^{\circ}) = -\frac{4}{\sqrt{2}+\sqrt{6}} = \sqrt{2} - \sqrt{6}$
\end{solution}
\end{problem}

\begin{problem} \label{sin105}
$\sin(105^{\circ})$

$\sin(105^{\circ}) = \answer{\frac{\sqrt{6}+\sqrt{2}}{4}}$
\end{problem}

\begin{problem}
$\csc(195^{\circ})$

\begin{solution}
$\csc(195^{\circ}) = \frac{4}{\sqrt{2}-\sqrt{6}} = -(\sqrt{2}+\sqrt{6})$
\end{solution}
\end{problem}

\begin{problem}
$\cot(255^{\circ})$

\begin{solution}
$\cot(255^{\circ}) = \frac{\sqrt{3}-1}{\sqrt{3}+1} = 2-\sqrt{3}$
\end{solution}
\end{problem}

\begin{problem}
$\tan(375^{\circ})$

\begin{solution}
$\tan(375^{\circ}) = \frac{3-\sqrt{3}}{3+\sqrt{3}} = 2-\sqrt{3}$
\end{solution}
\end{problem}

\begin{problem}
$\cos\left(\frac{13\pi}{12}\right)$

$\cos\left(\frac{13\pi}{12}\right) = \answer{-\frac{\sqrt{6}+\sqrt{2}}{4}}$
\end{problem}

\begin{problem}
$\sin\left(\frac{11\pi}{12}\right)$

\begin{solution}
$\sin\left(\frac{11\pi}{12}\right) = \frac{\sqrt{6} - \sqrt{2}}{4}$
\end{solution}
\end{problem}

\begin{problem}
$\tan\left(\frac{13\pi}{12}\right)$

\begin{solution}
$\tan\left(\frac{13\pi}{12}\right) = \frac{3-\sqrt{3}}{3+\sqrt{3}} = 2-\sqrt{3}$
\end{solution}
\end{problem}

\begin{problem}\label{cos7pi12} 
$\cos \left( \frac{7\pi}{12} \right)$

\begin{solution}
$\cos \left( \frac{7\pi}{12} \right) = \frac{\sqrt{2} - \sqrt{6}}{4}$
\end{solution}
\end{problem}

\begin{problem}
$\tan \left( \frac{17\pi}{12} \right)$

$\tan \left( \frac{17\pi}{12} \right) = \answer{2 + \sqrt{3}}$
\end{problem}

\begin{problem}\label{sinpi12}  
$\sin \left( \frac{\pi}{12} \right)$

\begin{solution}
$\sin \left( \frac{\pi}{12} \right) = \frac{\sqrt{6} - \sqrt{2}}{4}$
\end{solution}
\end{problem}

\begin{problem}
$\cot \left( \frac{11\pi}{12} \right)$

\begin{solution}
$\cot \left( \frac{11\pi}{12} \right) = -(2 + \sqrt{3})$
\end{solution}
\end{problem}

\begin{problem}
$\csc \left( \frac{5\pi}{12} \right)$

\begin{solution}
$\csc \left( \frac{5\pi}{12} \right) = \sqrt{6} - \sqrt{2}$
\end{solution}
\end{problem}

\begin{problem}\label{sumdifflast}
$\sec \left( -\frac{\pi}{12} \right)$ 

$\sec \left( -\frac{\pi}{12} \right) = \answer{\sqrt{6} - \sqrt{2}}$
\end{problem}


\end{question}

\begin{problem}
 If $\alpha$ is a Quadrant IV angle with $\cos(\alpha) = \frac{\sqrt{5}}{5}$, and  $\sin(\beta) = \frac{\sqrt{10}}{10}$, where $\frac{\pi}{2} < \beta < \pi$, find

\begin{enumerate}

\item $\cos(\alpha + \beta)$

\begin{solution}
$\cos(\alpha + \beta) = -\frac{\sqrt{2}}{10}$
\end{solution}

\item $\sin(\alpha + \beta)$

\begin{solution}
$\sin(\alpha + \beta) = \frac{7\sqrt{2}}{10}$
\end{solution}

\item $\tan(\alpha + \beta)$

\begin{solution}
$\tan(\alpha + \beta) = -7$ 
\end{solution}

\item $\cos(\alpha - \beta)$

\begin{solution}
$\cos(\alpha - \beta)= -\frac{\sqrt{2}}{2}$
\end{solution}

\item $\sin(\alpha - \beta)$

\begin{solution}
$\sin(\alpha - \beta) = \frac{\sqrt{2}}{2}$
\end{solution}

\item $\tan(\alpha - \beta)$

\begin{solution}
$\tan(\alpha - \beta) = -1$
\end{solution}

\end{enumerate}
\end{problem}

\begin{problem}
If $\csc(\alpha) = 3$, where $0 < \alpha < \frac{\pi}{2}$, and $\beta$ is a Quadrant II angle with $\tan(\beta) = -7$, find

\begin{enumerate}

\item $\cos(\alpha + \beta)$

\begin{solution}
$\cos(\alpha + \beta) = - \frac{4+7\sqrt{2}}{30}$
\end{solution}

\item $\sin(\alpha + \beta)$

\begin{solution}
$\sin(\alpha + \beta) = \frac{28-\sqrt{2}}{30}$
\end{solution}

\item $\tan(\alpha + \beta)$

\begin{solution}
$\tan(\alpha + \beta) = \frac{-28+\sqrt{2}}{4+7\sqrt{2}} = \frac{63-100\sqrt{2}}{41}$
\end{solution}

\item $\cos(\alpha - \beta)$

\begin{solution}
$\cos(\alpha - \beta) =  \frac{-4+7\sqrt{2}}{30}$
\end{solution}

\item $\sin(\alpha - \beta)$

\begin{solution}
$\sin(\alpha - \beta) = - \frac{28+\sqrt{2}}{30}$
\end{solution}

\item $\tan(\alpha - \beta)$

\begin{solution}
$\tan(\alpha - \beta)= \frac{28+\sqrt{2}}{4-7\sqrt{2}} = -\frac{63+100\sqrt{2}}{41}$
\end{solution}

\end{enumerate}
\end{problem}

\begin{problem}
If $\sin(\alpha) = \frac{3}{5}$, where $0 < \alpha < \frac{\pi}{2}$, and $\cos(\beta) = \frac{12}{13}$ where $\frac{3\pi}{2} < \beta < 2\pi$, find 

\begin{enumerate}

\item $\sin(\alpha + \beta)$

$\sin(\alpha + \beta) = \answer{\frac{16}{65}}$

\item $\cos(\alpha - \beta)$

$\cos(\alpha - \beta) = \answer{\frac{33}{65}}$

\item $\tan(\alpha - \beta)$

$\tan(\alpha - \beta) = \answer{\frac{56}{33}}$

\end{enumerate}
\end{problem}

\begin{problem}
If $\sec(\alpha) = -\frac{5}{3}$, where $\frac{\pi}{2} < \alpha < \pi$, and $\tan(\beta) = \frac{24}{7}$, where $\pi < \beta < \frac{3\pi}{2}$, find

\begin{enumerate}

\item $\csc(\alpha - \beta)$

\begin{solution}
$\csc(\alpha - \beta) = -\frac{5}{4}$
\end{solution}

\item $\sec(\alpha + \beta)$

\begin{solution}
$\sec(\alpha + \beta) = \frac{125}{117}$
\end{solution}

\item $\cot(\alpha + \beta)$

\begin{solution}
$\cot(\alpha + \beta) = \frac{117}{44}$
\end{solution}
\end{enumerate}
\end{problem}


\begin{question}
    
In Exercises \ref{expandedsinusoidexerfirst} - \ref{expandedsinusoidexerlast}, use Example \ref{expandedsinusoidex1} as a guide to show that the function is a sinusoid by rewriting it in the forms $C(t) = A \cos(\omega t + \phi) + B$ and $S(t) = A \sin(\omega t + \phi) + B$ for $\omega > 0$ and $0 \leq \phi < 2\pi$.

\begin{problem}\label{expandedsinusoidexerfirst}
$f(t) = \sqrt{2}\sin(t) + \sqrt{2}\cos(t) + 1$

\begin{solution}
$f(t) = \sqrt{2}\sin(t) + \sqrt{2}\cos(t) + 1 = 2\sin\left(t + \frac{\pi}{4}\right) + 1 = 2\cos\left(t + \frac{7\pi}{4}\right) + 1$ 
\end{solution}
\end{problem}

\begin{problem}
$f(t) = 3\sqrt{3}\sin(3t) - 3\cos(3t)$

\begin{solution}
$f(t) = 3\sqrt{3}\sin(3t) - 3\cos(3t) = 6\sin\left(3t + \frac{11\pi}{6}\right) = 6\cos\left(3t + \frac{4\pi}{3}\right)$
\end{solution}
\end{problem}

\begin{problem}
$f(t) = -\sin(t) + \cos(t) - 2$

\begin{solution}
$f(t) = -\sin(t) + \cos(t) - 2 = \sqrt{2}\sin\left(t + \frac{3\pi}{4}\right) - 2 = \sqrt{2}\cos\left(t + \frac{\pi}{4}\right) - 2$
\end{solution}
\end{problem}

\begin{problem}
$f(t) = -\frac{1}{2}\sin(2t) - \frac{\sqrt{3}}{2}\cos(2t)$

\begin{solution}
$f(t) = -\frac{1}{2}\sin(2t) - \frac{\sqrt{3}}{2}\cos(2t) = \sin\left(2t + \frac{4\pi}{3}\right) = \cos\left(2t + \frac{5\pi}{6}\right)$
\end{solution}
\end{problem}

\begin{problem}
$f(t) = 2\sqrt{3} \cos(t) - 2\sin(t)$

\begin{solution}
$f(t) = 2\sqrt{3} \cos(t) - 2\sin(t) = 4\sin\left(t+\frac{2\pi}{3}  \right)  = 4\cos\left(t + \frac{\pi}{6}\right)$
\end{solution}
\end{problem}

\begin{problem}
$f(t) = \frac{3}{2} \cos(2t) - \frac{3\sqrt{3}}{2} \sin(2t) + 6$

\begin{solution}
$f(t) = \frac{3}{2} \cos(2t) - \frac{3\sqrt{3}}{2} \sin(2t) + 6 =3\sin\left(2t + \frac{5\pi}{6}\right) + 6   = 3\cos\left(2t + \frac{\pi}{3}\right) + 6$
\end{solution}
\end{problem}

\begin{problem}
$f(t) = -\frac{1}{2} \cos(5t) -\frac{\sqrt{3}}{2} \sin(5t)$

\begin{solution}
$f(t) = -\frac{1}{2} \cos(5t) -\frac{\sqrt{3}}{2} \sin(5t) =  \sin\left(5t + \frac{7\pi}{6}\right) = \cos\left(5t + \frac{2\pi}{3}\right)$
\end{solution}
\end{problem}

\begin{problem}
$f(t) = -6\sqrt{3} \cos(3t) - 6\sin(3t) - 3$

\begin{solution}
$f(t) = -6\sqrt{3} \cos(3t) - 6\sin(3t) - 3 = 12\sin\left(3t + \frac{4\pi}{3}\right) - 3 = 12\cos\left(3t + \frac{5\pi}{6}\right) - 3$
\end{solution}
\end{problem}

\begin{problem}
$f(t) =  \frac{5\sqrt{2}}{2} \sin(t)  -\frac{5\sqrt{2}}{2} \cos(t)$

\begin{solution}
$f(t) =  \frac{5\sqrt{2}}{2} \sin(t)  -\frac{5\sqrt{2}}{2} \cos(t) = 5\sin\left(t + \frac{7\pi}{4}\right)= 5\cos\left(t + \frac{5\pi}{4}\right)$
\end{solution}
\end{problem}

\begin{problem}\label{expandedsinusoidexerlast}
$f(t) =3 \sin \left(\frac{t}{6}\right) -3\sqrt{3} \cos \left(\frac{t}{6}\right)$

\begin{solution}
$f(t) =3\sin\left(\frac{t}{6}\right) -3\sqrt{3} \cos\left(\frac{t}{6}\right) = 6\sin\left( \frac{t}{6}+\frac{5\pi}{3}\right)= 6\cos\left( \frac{t}{6}+\frac{7\pi}{6}\right) $
\end{solution}
\end{problem}
\end{question}

\begin{problem}\label{sinusoidexercise1}
    In Exercises \ref{expandedsinusoidexerfirst} - \ref{expandedsinusoidexerlast}, you should have noticed a relationship between the phases $\phi$ for the $S(t)$ and $C(t)$.  Show that if $f(t) = A \sin(\omega t + \alpha) + B$, then $f(t) = A \cos(\omega t + \beta) + B$ where $\beta = \alpha - \frac{\pi}{2}$. 
\end{problem}

\begin{problem}\label{sinusoidexercise2}
    Let $\phi$ be an angle measured in radians and let $P(a,b)$ be a point on the terminal side of $\phi$ when it is drawn in standard position.  Use Theorem \ref{cosinesinecircle} and the sum identity for sine in Theorem \ref{sinesumdifference} to show that  $f(t) = a \, \sin(\omega t) + b\, \cos(\omega t) + B$ (with  $\omega > 0$) can be rewritten as $f(t) = \sqrt{a^{2} + b^{2}}\sin(\omega t + \phi) + B$.
\end{problem}

\begin{problem}\label{twodaylightfunctions}
    In Example \ref{sinusoidsunlight} in Section \ref{GraphsofSineandCosine}, we developed two (seemingly) different formulas to model the hours of daylight, $H(t)$:   $H_{1}(t) = 9.25 \sin\left(\frac{\pi}{6} t - \frac{\pi}{2}\right) + 12.55$ and  $H_{2}(t) = -8.13 \sin\left(\frac{\pi}{6} t - 4.70\right)+ 12.5$.  Use the difference identities for sine to expand $H_{1}(t)$ and $H_{2}(t)$.  How different are they?
\end{problem}

\begin{question}
    
In Exercises \ref{identfirstident} - \ref{identlastident}, verify the identity.\footnote{Note: numbers \ref{conversionsineshift} and \ref{conversioncosineshift} are the conversion formulas stated in Theorem \ref{cosinesinefunctionprops} in Section \ref{GraphsofSineandCosine}.}

\begin{problem}\label{identfirstident} \label{conversionsineshift}
$\sin\left(\theta + \frac{\pi}{2}\right) = \cos(t)$  
\end{problem}

\begin{problem}\label{conversioncosineshift}
$\cos\left(\theta - \frac{\pi}{2} \right) = \sin(t)$    
\end{problem}

\begin{problem}
$\cos(\theta - \pi) = -\cos(\theta)$    
\end{problem}
 
\begin{problem}
$\sin(\pi - \theta) = \sin(\theta)$
\end{problem}

\begin{problem}
$\tan\left(\theta + \frac{\pi}{2} \right) = -\cot(\theta)$
\end{problem}

\begin{problem}
$\sin(\alpha + \beta) + \sin(\alpha - \beta) = 2\sin(\alpha)\cos(\beta)$  
\end{problem}

\begin{problem}
$\sin(\alpha + \beta) - \sin(\alpha - \beta) = 2\cos(\alpha) \sin(\beta)$
\end{problem}
 
\begin{problem}
$\cos(\alpha + \beta) + \cos(\alpha - \beta) = 2\cos(\alpha) \cos(\beta)$
\end{problem}

\begin{problem}
$\cos(\alpha + \beta) - \cos(\alpha - \beta) = -2\sin(\alpha) \sin(\beta)$
\end{problem}

\begin{problem}
$\frac{\sin(\alpha+\beta)}{\sin(\alpha-\beta)} = \frac{1+\cot(\alpha) \tan(\beta)}{1 - \cot(\alpha) \tan(\beta)}$    
\end{problem} 

\begin{problem}
$\frac{\cos(\alpha + \beta)}{\cos(\alpha - \beta)} = \frac{1 - \tan(\alpha)\tan(\beta)}{1 + \tan(\alpha)\tan(\beta)}$
\end{problem}

\begin{problem}\label{identlastident}
$\frac{\tan(\alpha + \beta)}{\tan(\alpha - \beta)} = \frac{\sin(\alpha)\cos(\alpha) + \sin(\beta)\cos(\beta)}{\sin(\alpha)\cos(\alpha) - \sin(\beta)\cos(\beta)}$
\end{problem}
\end{question}

\begin{question}
    
In Exercise \ref{derivativesine} - \ref{derivativecosine}, use the results from  Exercise \ref{sintovertexercise3} in Section \ref{TheOtherCircularFunctions}: $\lim_{\theta \rightarrow 0}$ $\frac{\sin(\theta)}{\theta} = 1$ and Exercise \ref{oneminuscosinelimit} in Section \ref{FundamentalTrigonometricIdentities} :   $\lim_{\theta \rightarrow 0}$ $\frac{1 - \cos(\theta)}{\theta} = 0$ to help you find the derivatives of the sine and cosine functions.

\begin{problem}\label{derivativesine}
\begin{enumerate}  
    \item  Verify for $f(t) = \sin(t)$, $\frac{f(t+h) - f(t)}{h} = \frac{\sin(t + h) - \sin(t)}{h} = \cos(t) \left(\frac{\sin(h)}{h} \right) + \sin(t) \left( \frac{\cos(h) - 1}{h} \right)$

\item  Show $f'(t) = \lim_{h \rightarrow 0} \frac{f(t+h) - f(t)}{h} = \cos(t)$.

\begin{hint}
$\lim_{\theta \rightarrow 0}$ $\frac{\cos(\theta) - 1}{\theta} =$ $\lim_{\theta \rightarrow 0}$ $\left(-\frac{1 - \cos(\theta)}{\theta}\right) = - (0) = 0$
\end{hint}

\item  Find the equation of the tangent line to the graph of $y = \sin(t)$ at $(0,0)$, $\left(\frac{\pi}{2}, 1\right)$ and $(\pi, 0)$.  

Check your answers graphically.

\begin{solution}
at $(0,0)$:  $y = x$;  at $\left(\frac{\pi}{2}, 1\right)$:  $y = 1$   at $(\pi, 0)$:  $y = -x + \pi$
\end{solution}

\end{enumerate}
\end{problem}

\begin{problem}\label{derivativecosine}

\begin{enumerate} 
\item   Verify for $g(t) = \cos(t)$, $\frac{g(t+h) - g(t)}{h} = \frac{\cos(t + h) - \cos(t)}{h} = \cos(t) \left( \frac{\cos(h) - 1}{h} \right) - \sin(t) \left(\frac{\sin(h)}{h} \right)$

\item  Show $g'(t) = \lim_{h \rightarrow 0} \frac{g(t+h) - g(t)}{h} =  - \sin(t)$.

\item  Find the equation of the tangent line to the graph of $y = \cos(t)$ at $(0,1)$, $\left(\frac{\pi}{2}, 0\right)$ and $\left(-\frac{\pi}{2}, 0\right)$.  

Check your answers graphically.

\begin{solution}
at $(0,1)$:  $y = 1$;  at $\left(\frac{\pi}{2}, 0\right)$:  $y = -x + \frac{\pi}{2}$   at $\left(-\frac{\pi}{2}, 0\right)$:  $y = x + \frac{\pi}{2}$
\end{solution}

\end{enumerate}
\end{problem}
\end{question}

\begin{problem}\label{derivativetangent}  

\begin{enumerate}  

\item\label{tanthetaoverthetalimit}  Use the fact\footnote{See Exercise \ref{sintovertexercise3} in Section \ref{TheOtherCircularFunctions}.} that $\lim_{\theta \rightarrow 0}$ $\frac{\sin(\theta)}{\theta} = 1$  to prove $\lim_{\theta \rightarrow 0}$ $\frac{\tan(\theta)}{\theta} = 1$.

\begin{hint}
$\frac{\tan(\theta)}{\theta} = \frac{\sin(\theta)}{\theta} \, \frac{1}{\cos(\theta)}$ \ldots
\end{hint} 

\item 
For $f(t) = \tan(t)$, show $\frac{f(t+h) - f(t)}{h} = \frac{\tan(t + h) - \tan(t)}{h} = \left( \frac{\tan(h)}{h} \right) \left(\frac{\sec^{2}(t)}{1 - \tan(t)\tan(h)} \right)$.
    
\item 
Use part \ref{tanthetaoverthetalimit} to show $f'(t) = \lim_{h \rightarrow 0} \frac{f(t+h) - f(t)}{h} = \sec^{2}(t)$ 

\item 
Find the equation of the tangent line to the graph of $y = \tan(t)$ at $(0,0)$, $\left(\frac{\pi}{4}, 1 \right)$ and $\left(-\frac{\pi}{4}, -1\right)$.  

Check your answers graphically.

\begin{solution}
at $(0,0)$:  $y = x$;  at $\left(\frac{\pi}{4}, 1\right)$:  $y = 2x - \frac{\pi}{2} +1$   at $\left(-\frac{\pi}{4}, -1\right)$:  $y = 2x + \frac{\pi}{2} - 1$
\end{solution}

\end{enumerate}
\end{problem}

\begin{question}
In Exercises \ref{idenhalfanglefirst} - \ref{idenhalfanglelast}, use the Half Angle Formulas to find the exact value.  You may have need of the Quotient, Reciprocal or Even / Odd Identities as well.

\begin{problem}\label{idenhalfanglefirst}
$\cos(75^{\circ})$  (compare with Exercise \ref{cos75})

$\cos(75^{\circ}) = \answer{\frac{\sqrt{2-\sqrt{3}}}{2}}$
\end{problem}

\begin{problem}
$\sin(105^{\circ})$  (compare with Exercise \ref{sin105})

\begin{solution}
$\sin(105^{\circ}) = \frac{\sqrt{2+\sqrt{3}}}{2}$
\end{solution}
\end{problem}

\begin{problem}
$\cos(67.5^{\circ})$

\begin{solution}
$\cos(67.5^{\circ})  = \frac{\sqrt{2-\sqrt{2}}}{2}$
\end{solution}
\end{problem}

\begin{problem}
$\sin(157.5^{\circ})$

\begin{solution}
$\sin(157.5^{\circ}) = \frac{\sqrt{2-\sqrt{2}}}{2}$
\end{solution}
\end{problem}

\begin{problem}
$\tan(112.5^{\circ})$

\begin{solution}
$\tan(112.5^{\circ}) = - \sqrt{\frac{2+\sqrt{2}}{2-\sqrt{2}}} = -1 - \sqrt{2}$
\end{solution}
\end{problem}

\begin{problem}
$\cos\left( \frac{7\pi}{12} \right)$  (compare with Exercise \ref{cos7pi12})

$\cos\left( \frac{7\pi}{12} \right) = \answer{-\frac{\sqrt{2-\sqrt{3}}}{2}}$  
\end{problem}

\begin{problem}
$\sin\left( \frac{\pi}{12} \right)$  (compare with Exercise \ref{sinpi12})

\begin{solution}
$\sin\left( \frac{\pi}{12} \right) = \frac{\sqrt{2-\sqrt{3}}}{2}$ 
\end{solution}
\end{problem}

\begin{problem}
$\cos \left( \frac{\pi}{8} \right)$

\begin{solution}
$\cos \left( \frac{\pi}{8} \right) = \frac{\sqrt{2 + \sqrt{2}}}{2}$
\end{solution}
\end{problem}

\begin{problem}
$\sin \left( \frac{5\pi}{8} \right)$

\begin{solution}
$\sin \left( \frac{5\pi}{8} \right) = \frac{\sqrt{2 + \sqrt{2}}}{2}$
\end{solution}
\end{problem}

\begin{problem}\label{idenhalfanglelast}
$\tan \left( \frac{7\pi}{8} \right)$ 

\begin{solution}
$\tan \left( \frac{7\pi}{8} \right) = -\sqrt{ \frac{2 - \sqrt{2}}{2 + \sqrt{2}} } =1-\sqrt{2}$
\end{solution}
\end{problem}

\end{question}

\begin{question}
    
In Exercises \ref{doublehalffirst} - \ref{doublehalflast}, use the given information about $\theta$ to find the exact values of 

\begin{itemize}

\item $\sin(2\theta)$
\item $\sin\left(\frac{\theta}{2}\right)$
\item $\cos(2\theta)$
\item $\cos\left(\frac{\theta}{2}\right)$
\item $\tan(2\theta)$
\item $\tan\left(\frac{\theta}{2}\right)$

\end{itemize}

\begin{problem}\label{doublehalffirst}
$\sin(\theta) = -\frac{7}{25}$ where $\frac{3\pi}{2} < \theta < 2\pi$ 

\begin{solution}
\begin{itemize}

\item $\sin(2\theta) = -\frac{336}{625}$
\item $\sin\left(\frac{\theta}{2}\right) = \frac{\sqrt{2}}{10}$
\item $\cos(2\theta) = \frac{527}{625}$
\item $\cos\left(\frac{\theta}{2}\right) = -\frac{7\sqrt{2}}{10}$
\item $\tan(2\theta) = -\frac{336}{527}$
\item $\tan\left(\frac{\theta}{2}\right) = -\frac{1}{7}$

\end{itemize}
\end{solution}
\end{problem}

\begin{problem}
$\cos(\theta) = \frac{28}{53}$ where $0 < \theta < \frac{\pi}{2}$

\begin{solution}
\begin{itemize}

\item $\sin(2\theta) = \frac{2520}{2809}$
\item $\sin\left(\frac{\theta}{2}\right) = \frac{5\sqrt{106}}{106}$
\item $\cos(2\theta) = -\frac{1241}{2809}$
\item $\cos\left(\frac{\theta}{2}\right) = \frac{9\sqrt{106}}{106}$
\item $\tan(2\theta) = -\frac{2520}{1241}$
\item $\tan\left(\frac{\theta}{2}\right) = \frac{5}{9}$

\end{itemize}
\end{solution}
\end{problem}

\begin{problem}
$\tan(\theta) = \frac{12}{5}$ where $\pi < \theta < \frac{3\pi}{2}$

\begin{itemize}

\item $\sin(2\theta) = \answer{\frac{120}{169}}$
\item $\sin\left(\frac{\theta}{2}\right) = \answer{\frac{3\sqrt{13}}{13}}$
\item $\cos(2\theta) = \answer{-\frac{119}{169}}$
\item $\cos\left(\frac{\theta}{2}\right) = \answer{-\frac{2\sqrt{13}}{13}}$
\item $\tan(2\theta) = \answer{-\frac{120}{119}}$
\item $\tan\left(\frac{\theta}{2}\right) = \answer{-\frac{3}{2}}$

\end{itemize}
\end{problem}

\begin{problem}
$\csc(\theta) = 4$ where $\frac{\pi}{2} < \theta < \pi$

\begin{solution}
\begin{itemize}

\item $\sin(2\theta) = -\frac{\sqrt{15}}{8}$
\item $\sin\left(\frac{\theta}{2}\right) =\frac{\sqrt{8+2\sqrt{15}}}{4} \\ \phantom{\tan\left(\frac{\theta}{2}\right) = 4+\sqrt{15}}$ 
\item $\cos(2\theta) = \frac{7}{8}$
\item $\cos\left(\frac{\theta}{2}\right) = \frac{\sqrt{8-2\sqrt{15}}}{4} \\ \phantom{\tan\left(\frac{\theta}{2}\right) = 4+\sqrt{15}} $
\item $\tan(2\theta) = -\frac{\sqrt{15}}{7}$
\item $\tan\left(\frac{\theta}{2}\right) = \sqrt{\frac{8+2\sqrt{15}}{8-2\sqrt{15}}} \\ \tan\left(\frac{\theta}{2}\right) = 4+\sqrt{15}$

\end{itemize}
\end{solution}
\end{problem}

\begin{problem}
$\cos(\theta) = \frac{3}{5}$ where $0 < \theta < \frac{\pi}{2}$

\begin{solution}
\begin{itemize}

\item $\sin(2\theta) = \frac{24}{25}$
\item $\sin\left(\frac{\theta}{2}\right) = \frac{\sqrt{5}}{5}$
\item $\cos(2\theta) = -\frac{7}{25}$
\item $\cos\left(\frac{\theta}{2}\right) = \frac{2\sqrt{5}}{5}$
\item $\tan(2\theta)=-\frac{24}{7} $
\item $\tan\left(\frac{\theta}{2}\right) = \frac{1}{2}$

\end{itemize}
\end{solution}
\end{problem}

\begin{problem}
$\sin(\theta) = -\frac{4}{5}$ where $\pi < \theta < \frac{3\pi}{2}$

\begin{solution}
\begin{itemize}

\item $\sin(2\theta) = \frac{24}{25}$
\item $\sin\left(\frac{\theta}{2}\right) = \frac{2\sqrt{5}}{5}$
\item $\cos(2\theta) = -\frac{7}{25}$
\item $\cos\left(\frac{\theta}{2}\right) = -\frac{\sqrt{5}}{5}$
\item $\tan(2\theta)=-\frac{24}{7} $
\item $\tan\left(\frac{\theta}{2}\right) = -2$

\end{itemize}
\end{solution}
\end{problem}

\begin{problem}
$\cos(\theta) = \frac{12}{13}$ where $\frac{3\pi}{2} < \theta < 2\pi$

\begin{solution}
\begin{itemize}

\item $\sin(2\theta) = -\frac{120}{169}$
\item $\sin\left(\frac{\theta}{2}\right) = \frac{\sqrt{26}}{26}$
\item $\cos(2\theta) = \frac{119}{169}$
\item $\cos\left(\frac{\theta}{2}\right) = -\frac{5\sqrt{26}}{26}$
\item $\tan(2\theta)=-\frac{120}{119}$
\item $\tan\left(\frac{\theta}{2}\right) = -\frac{1}{5}$

\end{itemize}
\end{solution}
\end{problem}

\begin{problem}
$\sin(\theta) = \frac{5}{13}$ where $\frac{\pi}{2} < \theta < \pi$

\begin{solution}
\begin{itemize}

\item $\sin(2\theta) = -\frac{120}{169}$
\item $\sin\left(\frac{\theta}{2}\right) = \frac{5\sqrt{26}}{26}$
\item $\cos(2\theta) = \frac{119}{169}$
\item $\cos\left(\frac{\theta}{2}\right) = \frac{\sqrt{26}}{26}$
\item $\tan(2\theta)=-\frac{120}{119}$
\item $\tan\left(\frac{\theta}{2}\right) = 5$

\end{itemize}
\end{solution}
\end{problem}

\begin{problem}
$\sec(\theta) = \sqrt{5}$ where $\frac{3\pi}{2} < \theta < 2\pi$

\begin{solution}
\begin{itemize}

\item $\sin(2\theta) = -\frac{4}{5}$
\item $\sin\left(\frac{\theta}{2}\right) = \frac{\sqrt{50-10\sqrt{5}}}{10} \\ \phantom{\tan\left(\frac{\theta}{2}\right) =\frac{5-5\sqrt{5}}{10}}$ 

\item $\cos(2\theta) = -\frac{3}{5}$
\item $\cos\left(\frac{\theta}{2}\right)= -\frac{\sqrt{50+10\sqrt{5}}}{10} \\ \phantom{\tan\left(\frac{\theta}{2}\right) =\frac{5-5\sqrt{5}}{10}}$ 
\item $\tan(2\theta)=\frac{4}{3}$
\item $\tan\left(\frac{\theta}{2}\right) =  -\sqrt{\frac{5-\sqrt{5}}{5+\sqrt{5}}} \\ \tan\left(\frac{\theta}{2}\right) =\frac{5-5\sqrt{5}}{10}$

\end{itemize}
\end{solution}
\end{problem}

\begin{problem}\label{doublehalflast}
$\tan(\theta) = -2$ where $\frac{\pi}{2} < \theta < \pi$

\begin{solution}
\begin{itemize}

\item $\sin(2\theta) = -\frac{4}{5}$
\item $\sin\left(\frac{\theta}{2}\right) = \frac{\sqrt{50+10\sqrt{5}}}{10} \\ \phantom{\tan\left(\frac{\theta}{2}\right) =\frac{5-5\sqrt{5}}{10}}$ 

\item $\cos(2\theta) = -\frac{3}{5}$
\item $\cos\left(\frac{\theta}{2}\right)= \frac{\sqrt{50-10\sqrt{5}}}{10} \\ \phantom{\tan\left(\frac{\theta}{2}\right) =\frac{5-5\sqrt{5}}{10}}$ 
\item $\tan(2\theta)=\frac{4}{3}$
\item $\tan\left(\frac{\theta}{2}\right) =  \sqrt{\frac{5+\sqrt{5}}{5-\sqrt{5}}} \\ \tan\left(\frac{\theta}{2}\right) =\frac{5+5\sqrt{5}}{10}$

\end{itemize}
\end{solution}
\end{problem}

\end{question}

\begin{question}
    
In Exercises \ref{moreidentfirst} - \ref{moreidentlast}, verify the identity.  Assume all quantities are defined.

\begin{problem}\label{moreidentfirst}
$(\cos(\theta) + \sin(\theta))^2 = 1 + \sin(2\theta)$
\end{problem}

\begin{problem}
    $(\cos(\theta) - \sin(\theta))^2 = 1 - \sin(2\theta)$
\end{problem}

\begin{problem}
    $\tan(2t) = \frac{1}{1-\tan(t)} - \frac{1}{1+\tan(t)}$
\end{problem}

\begin{problem}
    $\csc(2\theta) = \frac{\cot(\theta) + \tan(\theta)}{2}$
\end{problem}

\begin{problem}
    $8 \sin^{4}(x) = \cos(4x) - 4\cos(2x)+3$
\end{problem}

\begin{problem}
    $8 \cos^{4}(x) = \cos(4x) + 4\cos(2x)+3$
\end{problem}

\begin{problem}\label{sine3theta}
    $\sin(3\theta) = 3\sin(\theta) - 4\sin^{3}(\theta)$
\end{problem}

\begin{problem}
    $\sin(4\theta) = 4\sin(\theta)\cos^{3}(\theta) - 4\sin^{3}(\theta)\cos(\theta)$
\end{problem}

\begin{problem}
    $32\sin^{2}(t) \cos^{4}(t) = 2 + \cos(2t) - 2\cos(4t) - \cos(6t)$
\end{problem}

\begin{problem}
    $32\sin^{4}(t) \cos^{2}(t) = 2 - \cos(2t) - 2\cos(4t) + \cos(6t)$
\end{problem}

\begin{problem}\label{cosine4theta} 
    $\cos(4\theta) = 8\cos^{4}(\theta) - 8\cos^{2}(\theta) + 1$
\end{problem}

\begin{problem}
    $\cos(8\theta) = 128\cos^{8}(\theta)-256\cos^{6}(\theta)+160\cos^{4}(\theta)-32\cos^{2}(\theta)+1$ (HINT:  Use the result to \ref{cosine4theta}.)
\end{problem}

\begin{problem}
    $\sec(2x) = \frac{\cos(x)}{\cos(x) + \sin(x)} + \frac{\sin(x)}{\cos(x)-\sin(x)}$ 
\end{problem}

\begin{problem}
     $\frac{1}{\cos(\theta) - \sin(\theta)} + \frac{1}{\cos(\theta) + \sin(\theta)} = \frac{2\cos(\theta)}{\cos(2\theta)}$
\end{problem}

\begin{problem} \label{moreidentlast}
    $\frac{1}{\cos(\theta) - \sin(\theta)} - \frac{1}{\cos(\theta) + \sin(\theta)} = \frac{2\sin(\theta)}{\cos(2\theta)}$
\end{problem}


\end{question}

\begin{problem}\label{preludetoarctrigsine}
Suppose $\theta$ is a Quadrant I angle with $\sin(\theta) = x$. Verify the following formulas

\begin{enumerate}

\item  $\cos(\theta) = \sqrt{1-x^2}$

\item  $\sin(2\theta) = 2x\sqrt{1-x^2}$

\item $\cos(2\theta) = 1 - 2x^2$

\end{enumerate}
\end{problem}

\begin{problem}
Discuss with your classmates how each of the formulas, if any, in Exercise \ref{preludetoarctrigsine} change if we change assume $\theta$ is a Quadrant II, III, or IV angle.
\end{problem}

\begin{problem}\label{preludetoarctrigtan}
Suppose $\theta$ is a Quadrant I angle with $\tan(\theta) = x$. Verify the following formulas
\begin{enumerate}

\item $\cos(\theta) = \frac{1}{\sqrt{x^2+1}}$
\item $\sin(\theta) = \frac{x}{\sqrt{x^2+1}}$
\item $\sin(2\theta) = \frac{2x}{x^2+1}$
\item $\cos(2\theta) = \frac{1-x^2}{x^2+1}$

\end{enumerate}

\end{problem}

\begin{problem}
Discuss with your classmates how each of the formulas, if any, in Exercise \ref{preludetoarctrigtan} change if we change assume $\theta$ is a Quadrant II, III, or IV angle.
\end{problem}

\begin{problem}
If $\sin(t) = x$ for  $-\frac{\pi}{2} < t < \frac{\pi}{2}$, find an expression for $\tan(t)$ in terms of $x$.

\begin{solution}
$\tan(t) = \frac{x}{\sqrt{1 - x^2}}$
\end{solution}
\end{problem}

\begin{problem}
If $\tan(\theta) = x$ for $-\frac{\pi}{2} < \theta < \frac{\pi}{2}$,  find an expression for $\sec(\theta)$ in terms of $x$.

\begin{solution}
$\sec(\theta) = \sqrt{1+x^2}$ 
\end{solution}
\end{problem}

\begin{problem}
If $\sec(\theta) = x$ where  $\theta$ is a Quadrant II angle, find  an expression for $\tan(\theta)$ in terms of $x$.

$\tan(\theta) = \answer{-\sqrt{x^2-1}}$
\end{problem}

\begin{problem}
If $\sin(t) = \frac{x}{2}$ for $-\frac{\pi}{2} < t < \frac{\pi}{2}$, find an expression for $\cos(2t)$ in terms of $x$.

\begin{solution}
$\cos(2t) = 1 - \frac{x^{2}}{2}$
\end{solution}
\end{problem}

\begin{problem}
If $\tan(\theta) = \frac{x}{7}$ for $-\frac{\pi}{2} < \theta < \frac{\pi}{2}$, find an expression for $\sin(2\theta)$ in terms of $x$.

\begin{solution}
$\sin(2\theta) = \frac{14x}{x^{2} + 49}$
\end{solution}
\end{problem}

\begin{problem}
If $\sec(t) = \frac{x}{4}$ for $0 < t < \frac{\pi}{2}$, find an expression for $\ln|\sec(t) + \tan(t)|$ in terms of $x$.

\begin{solution}
$\ln|\sec(t) + \tan(t)| = \ln |x + \sqrt{x^{2} + 16}| - \ln(4)$
\end{solution}
\end{problem}

\begin{problem}
Show that $\cos^{2}(\theta) - \sin^{2}(\theta) = 2\cos^{2}(\theta) - 1 = 1 - 2\sin^{2}(\theta)$ for all $\theta$.
\end{problem}

\begin{problem}
Let $\theta$ be a Quadrant III angle with $\cos(\theta) = -\frac{1}{5}$.  Show that this is not enough information to determine the sign of  $\sin\left(\frac{\theta}{2}\right)$ by first assuming $3\pi < \theta < \frac{7\pi}{2}$ and then assuming $\pi < \theta < \frac{3\pi}{2}$ and computing $\sin\left(\frac{\theta}{2}\right)$ in both cases.
\end{problem}

\begin{problem}
Without using your calculator, show that $\frac{\sqrt{2 + \sqrt{3}}}{2} = \frac{\sqrt{6} + \sqrt{2}}{4}$
\end{problem}

\begin{problem}
In part \ref{cosinepolynomial} of Example \ref{doubleangleex}, we wrote $\cos(3\theta)$ as a polynomial in terms of $\cos(\theta)$.  In Exercise \ref{cosine4theta}, we had you verify an identity which expresses $\cos(4\theta)$ as a polynomial in terms of $\cos(\theta)$.   Can you find a polynomial in terms of $\cos(\theta)$ for $\cos(5\theta)$?  $\cos(6\theta)$?  Can you find a pattern so that $\cos(n\theta)$ could be written as a polynomial in cosine for any natural number $n$?
\end{problem}

\begin{problem}
In Exercise \ref{sine3theta}, we has you verify an identity which expresses $\sin(3\theta)$ as a polynomial in terms of $\sin(\theta)$.   Can you do the same for  $\sin(5\theta)$?  What about for $\sin(4\theta)$?  If not, what goes wrong?
\end{problem}


\begin{question}
In Exercises \ref{idengraphfirst} - \ref{idengraphlast}, verify the identity by graphing the right and left hand using a graphing utility.

\begin{problem}\label{idengraphfirst} 
$\sin^{2}(t) + \cos^{2}(t) = 1$
\end{problem}

\begin{problem}
$\sec^{2}(x) - \tan^{2}(x) = 1$
\end{problem}

\begin{problem}
$\cos(t) = \sin\left(\frac{\pi}{2} - t\right)$
\end{problem}

\begin{problem}
$\tan(x+\pi) = \tan(x)$ 
\end{problem}
 
\begin{problem}
$\sin(2t) = 2\sin(t)\cos(t)$
\end{problem}  

\begin{problem}\label{idengraphlast}
$\tan\left(\frac{x}{2}\right) = \frac{\sin(x)}{1+\cos(x)}$ 
\end{problem}

\end{question}

\begin{question}
In Exercises \ref{prodsumfirst} - \ref{prodsumlast}, write the  given product as a sum. Note: you may need to use an Even/Odd Identity to match the answer provided.

\begin{problem}\label{prodsumfirst}
$\cos(3\theta)\cos(5\theta)$ 

\begin{solution}
$\frac{\cos(2\theta) + \cos(8\theta)}{2}$
\end{solution}
\end{problem}

\begin{problem}
$\sin(2t)\sin(7t)$

\begin{solution}
$\frac{\cos(5t) - \cos(9t)}{2}$
\end{solution}
\end{problem}

\begin{problem}
$\sin(9x)\cos(x)$

$\answer{\frac{\sin(8x) + \sin(10x)}{2}}$
\end{problem}

\begin{problem}
$\cos(2\theta) \cos(6\theta)$

\begin{solution}
$\frac{\cos(4\theta) + \cos(8\theta)}{2}$
\end{solution}
\end{problem}

\begin{problem}
$\sin(3t) \sin(2t)$

\begin{solution}
$\frac{\cos(t) - \cos(5t)}{2}$
\end{solution}
\end{problem}

\begin{problem}\label{prodsumlast}
$\cos(x) \sin(3x)$ 

\begin{solution}
$\frac{\sin(2x) + \sin(4x)}{2}$
\end{solution}
\end{problem}

\end{question}

\begin{question}
In Exercises \ref{sumprodfirst} - \ref{sumprodlast},  write the given sum as a product. Note:  you may need to use an Even/Odd or Cofunction Identity to match the answer provided.

\begin{problem}\label{sumprodfirst}
$\cos(3\theta) + \cos(5\theta)$

\begin{solution}
$2\cos(4\theta)\cos(\theta)$
\end{solution}
\end{problem}

\begin{problem}
$\sin(2t) - \sin(7t)$

\begin{solution}
$-2\cos \left( \frac{9}{2}t \right) \sin \left( \frac{5}{2}t \right)$
\end{solution}
\end{problem}

\begin{problem}
$\cos(5x) - \cos(6x)$

\begin{solution}
$2\sin \left( \frac{11}{2}x \right) \sin \left( \frac{1}{2}x \right)$
\end{solution}
\end{problem}

\begin{problem}
$\sin(9x) - \sin(-x)$

$\answer{2\cos(4x)\sin(5x)}$
\end{problem}

\begin{problem}
$\sin(t) + \cos(t)$

\begin{solution}
$\sqrt{2}\cos \left(t - \frac{\pi}{4} \right)$
\end{solution}
\end{problem}

\begin{problem}\label{sumprodlast}
$\cos(x) - \sin(x)$

\begin{solution}
$-\sqrt{2}\sin \left(x - \frac{\pi}{4} \right)$
\end{solution}
\end{problem}

\end{question}

\begin{question}
In Exercises \ref{beatexfirst} - \ref{beatexlast}, using the remarks following Example \ref{expandedsinusoidex1} on page \pageref{beats} as a guide,  rewrite the given function $f(t)$ as a product of sinusoids. Identify the functions which create the `wave envelope.' Check your answer by graphing  the function along with the `wave-envelope' using a graphing utility.


\begin{problem}\label{beatexfirst}
$f(t) = \cos(3t) + \cos(5t)$

\begin{solution}
$f(t) = [2\cos(t)] \cos(4t)$, $A(t) = 2\cos(t)$, wave-envelope:  $y = \pm 2\cos(t)$
\end{solution}
\end{problem}  

\begin{problem}
$f(t) = 3\cos(5t) - 3\cos(6t)$

\begin{solution}
$f(t) = \left[6\sin \left( \frac{1}{2} t \right) \right] \sin \left( \frac{11}{2} t \right) $, $A(t) = 6\sin \left( \frac{1}{2} t \right) $, wave-envelope:  $y = \pm  6\sin \left( \frac{1}{2} t \right)$
\end{solution}
\end{problem}

\begin{problem}
$f(t) = \frac{1}{2} \sin(9t) + \frac{1}{2} \sin(t)$

\begin{solution}
$f(t) = [\cos(4t)] \sin(5t)$, $A(t) = \cos(4t)$, wave-envelope:  $y = \pm \cos(4t)$
\end{solution}
\end{problem}

\begin{problem}\label{beatexlast}
$f(t) = \frac{2}{3}\sin(2t) - \frac{2}{3}\sin(7t)$

\begin{solution}
$f(t) =  \left[-\frac{4}{3}\sin \left( \frac{5}{2} t \right) \right]\cos \left( \frac{9}{2} t \right) $, $A(t) = -\frac{4}{3}\sin \left( \frac{5}{2} t \right)$, wave-envelope:  $y = \pm \frac{4}{3} \sin \left( \frac{5}{2} t \right)$
\end{solution}
\end{problem}

\end{question}

\begin{problem}
Verify the Even / Odd Identities for tangent, secant, cosecant and cotangent.
\end{problem}
 
\begin{problem}
Verify the Cofunction Identities for tangent, secant, cosecant and cotangent.
\end{problem}

\begin{problem}
Verify the Difference Identities for sine and tangent.
\end{problem}

\begin{problem}
Verify the Product to Sum Identities.
\end{problem}

\begin{problem}
Verify the Sum to Product Identities.
\end{problem}

\end{document}